\usection{Il Vuoto Nichilistico}

\dnd{} \wargam{oso} ha una caratteristica sfacciata che è piuttosto rara da incontrare nei giochi di ruolo: prevede vittoria e sconfitta reali. È reale nello stesso senso in cui le competizioni sportive sono reali: il successo dona gloria e la sconfitta causa vergogna e frustrazione. Perché il gioco sia potente e significativo, trasformativo, devi entrare in questa mentalità, deve importartene fino a questo punto.

Che il gioco sia esaltante nel successo e deludente nel fallimento è, credo, al centro del più grande paradosso e pericolo che lo stile del \wargam{e} si trova davanti: se vincere è davvero divertente, perché non continuare a farlo? Che razza di uomo è l'arbitro che gioisce nella sconfitta? Penso che \dnd{} abbia storicamente abbandonato le sfide reali perché ci vogliono sforzo e chiarezza mentale per mantenere l'ideale paradossale della posizione neutrale dell'arbitro: un funzionario del gioco che \textit{ovviamente} vorrebbe vedere i giocatori avere successo e prosperare, ma che \textit{nonostante questo} osserva il procedere del gioco come un falco, abbattendo senza remore ciò che è debole, sfortunato e falso.

\usubsection{Graziosità nella sconfitta}

Come chiunque giochi davvero a \dnd{} seguendone le regole, ho avuto la mia parte di sconfitte e fallimenti nel corso degli anni. Alcune volte si è trattato di eccitanti corse verso la lina di arrivo, cadendo solo all'ultimo momento; altre volte erano buffe commedie di errori. E altre volte ancora si è trattato di inciampi completamente privi di emozione e significato, in cui personaggi promettenti, spunti di avventura eccitanti e posizioni di campagna uniche sono stati annichiliti senza alcuna buona ragione. I dadi sono semplicemente andati in quella direzione ed eccoci qua, persi oltre l'umana ragionevolezza. Un sacco di sforzo e di speranze sono entrati nel dungeon e l'unica cosa che ne è uscita è stato\dots il nulla. Non è riemerso niente.

Il mio nomignolo per la natura inumana della sconfitta in \dnd{} è \textbf{il vuoto nichilistico}. È un mostro di una natura leggermente differente da quelli che incontri nella maggior parte degli sport, poiché \dnd{} non si gioca contro partecipanti umani e quindi non c'è nessuno a celebrare una vittoria quando la nostra squadra perde. C'è solo il vuoto. È persino peggio di molti sport solitari come l'atletica leggera o il golf e simili, poiché la preparazione è così ardua e ogni successo di scopo maggiore si costruisce necessariamente su successi precedenti. Le comparazioni migliori sono con hobby come l'arrampicata o la missilistica amatoriale, dove il fallimento nel raggiungere la cima significa sprecare tutto l'allenamento e la preparazione che si sono messi nella spedizione. Rischi il vuoto nichilistico innalzandoti verso la gloria del successo.

Talvolta, mi si sente descrivere il gioco come un masochistico esercizio di fallimento, affrontato da coloro a cui piace la liquirizia salata e, generalmente, si sentono vivi solo quando sono in difficoltà. Incoraggio assolutamente la gente a provare, ma se quella prima sconfitta non fa per te, se ti lascia semplicemente depresso, allora magari è meglio se giochiamo ad altro, ci sono un sacco di giochi là fuori che non trattano la moneta del vuoto.

Ma per coloro che possono affrontarla, c'è il risvolto positivo, la gloriosa verità: ogni volta che perdiamo, l'incontro col vuoto afferma l'autenticità del nostro gioco. Il sacramento del vuoto, quando lo incontriamo, siamo noi che tocchiamo la realtà del gioco: è \textit{davvero} un gioco di vittoria e sconfitta e possiamo davvero fare le mosse sbagliate e perdere in questo esercizio. Non è teatro, l'arbitro non bara per noi. E questo significa che quando vinciamo, anche la vittoria è reale.

\usubsection{Non ti abbandonerò}

La perspicace teorica dei GDR Meguey Baker una volta ha classificato i vari contratti sociali tra gruppi di gioco in due categorie:

\textbf{Nessuno si fa male} è un gioco in cui le poste emotive in gioco sono basse, ci sono degli efficaci guard rail tecnici e i colpi sono trattenuti. Se ferisci i sentimenti di un altro giocatore, è colpa tua e dovresti chiedere scusa e aggiustare il tiro.

\textbf{Non ti abbandonerò} è un gioco che, in un certo senso, punta a far male. I giocatori, consensualmente, si impegnano emotivamente e giocano a fondo ed è possibile sentirsi feriti. Quando succede va bene, ce lo aspettavamo; il gruppo non lo renderà strano né si disassocierà dal dolore, ci si lavora su assieme.

In alcuni posti, l'analisi di Meg è stata interpretata nel senso che c'è un buon tipo di gioco e un cattivo tipo di gioco, ma questo \textit{ovviamente} non è lo spirito dell'analisi. Mi permetto di dirlo perché il tipo di gioco che sto cercando di descrivere qui appartiene alla seconda categoria: l'approccio da \wargam{e}, quando affrontato come una potente attività con il potenziale per la crescita personale, potrebbe tranquillamente ferirti. È probabilmente uno dei motivi principali per cui noi umani gareggiamo negli sport, allenarsi e imparare a sopportare il fallimento, tirarci su e continuare a vivere. Non si può fare se nessuno si fa male.

Questo non vuol dire che non si possa giocare "nessuno si fa male" in maniera \wargam{e}. Ovviamente si può. Basta assicurarsi che l'argomento e i processi del gioco siano costruiti in maniera da alienare emotivamente i giocatori dagli eventi di gioco, tanto la vittoria quanto la sconfitta. Mantenere un approccio casual. Magari eliminare la morte del personaggio: è uno specifico dettagli d'alto profilo che spesso risulta un punto d'inciampo in questo senso. O mantenere le relazioni coi personaggi su un piano scherzoso, senza prestare troppa attenzione a ciò che gli accade.

Ma è fondamentalmente una scelta tra l'essere casual e fare sul serio. Per ottenere il massimo da questo gioco, o qualsiasi gioco, a un certo punto bisogna spingersi oltre. Per farlo, credo che Meg abbia assolutamente ragione riguardo al giusto tipo di contratto sociale da adottare: i giocatori devono supportarsi a vicenda ed esserci gli uni per gli altri se qualcuno si fa male.

Cosa voglia dire "non abbandonare" cambia moltissimo tra culture diverse, quindi non ti dirò di andare ad abbracciare un giocatore che ha appena perso un personaggio amato; non è detto che il tuo gruppo sia così. I miei, di solito, non lo sono; i finlandesi non sono esattamente affettuosi. Qua di solito si tratta più di lasciare che l'altro giocatore sobbolisca per il resto della sessione senza attirare l'attenzione sul fatto che è ferito e vedere dopo se vuole parlarne. Tutto molto situazionale.

Posso, comunque, descrivere come si presenti spesso l'abbandono emotivo: quando sei a disagio con il dolore di un altro giocatore, lo sminuisci in un tentativo di forzare la partita a tornare casual. Magari pensi che questo sia come si gioca "nessuno si farà male" e poi sei un mezzo stronzo che pensa che è colpa del giocatore ferito se si sente ferito, quindi ora tocca a te fare uso di un metaforico bastone per fargli realizzare dov'è e smettere di piangere come una bambina. Quindi scherzi sulla sua perdita, magari dicendo qualcosa come "non è nulla di che, sii uomo" o simili. Sminuisci il suo dolore.

Sicuramente funziona nel senso che quel giocatore, in futuro, eviterà accuratamente di impegnarsi emotivamente col gioco. Impararerà a giocare casual perché saprà che, quando gioca con te, la ricompensa per prendere le cose sul serio e magari rimanerci un po' male per una sconfitta è l'umiliazione.

Tutto questo vale anche per cose diverse dal perdere un personaggio in \dnd{}, a proposito. Se stai pensando a situazioni in cui i giocatori si sentono feriti e perdono la loro compostezza, queste sono esattamente le situazioni di cui stiamo parlando. Ma, ovviamente, se c'è un aspetto regolare del gioco che può prenderti alla sprovvista e frustrare anche le anime più insensibili, è proprio il vuoto nichilistico della sconfitta.

\usubsection{Il punto critico di mid-tier}

Il vuoto nichilistico ha un grande effetto sulla natura del gioco; una volta che lo affronti e lo vedi, l'intero gioco non può fare a meno di orientarsi verso di esso. È la natura del gioco, quindi non c'è nulla di cui preoccuparsi. Imparerete ad apprezzare l'autenticità del prendersi un metaforico calcio nelle palle, oppure troverete un gioco più delicato da giocare al suo posto.

Tuttavia, vorrei anche affermare che esiste un paradosso creativo insito nella natura di \dnd{} come \wargam{e}, un'insidiosa interazione emergente che potrebbe persino spiegare molto di come il gioco sia cambiato storicamente. Mi piace chiamare questa caratteristica emergente il \textbf{punto critico di mid-tier} per motivi che saranno presto evidenti.

Il problema è che la posta emotiva in gioco non è costante; più un avventuriero progredisce, più sforzo vi viene investito dal giocatore. Per questa ragione, nessuno può semplicemente dire di essersi immunizzato contro il vuoto nichilistico. Magari ti sta bene perdere un personaggio appena creato; magari ti sta bene perdere un personaggio che ha incrociato qualche successo minore. Ma non possiamo mai dire che sei immune al vuoto nella sua interezza perché il gioco non sembra avere alcun limite effettivo a quanto a lungo puoi cadere: più in alto sali, più ti coinvolgi emotivamente, più forte è la caduta.

Contemporaneamente, l'unica costante del gioco è la sua condizione di sconfitta: perdere il personaggio. Questa è la condizione di base contro cui si misura il successo. Da qui il paradosso creativo: per avere più successo in \dnd{}, devi accettare di giocare con una posta sempre più alta, con relativamente meno successo da ottenere, contro un rischio di sconfitta sempre maggiore.

Dipende dai dettagli, ma penso che a un certo punto della campagna si arrivi a un momento in cui la domanda paradossale di sempre più rischi per sempre meno successi diventi inaccettabile. Penso di averla vista incompere, genericamente, da qualche parte ai livelli medi, tra i 10.000 e i 100.000 PX totali. Hai messo così tanto lavoro in un personaggio che rischiarne la vita per ulteriori guadagni marginali inizia a sembrare folle per un'osservatore esterno, non preso dalla fame dell'azzardo.

Credo che l'esistenza di questa caratteristica emergente del gioco spieghi un bel po' del perché i testi di gioco siano come sono: cercano di evitare il vuoto nichilistico e lasciarti lo spazio per continuare a giocare. Se sia o meno qualcosa da fare è, ovviamente, una domanda altamente filosofica. Ci si addentra nella natura del gioco.

Qua ci sono le diverse idee che ho visto per come risolvere questo punto di crisi da mid-tier:

\paragraph{Il gioco ha un limite di livello}: potrebbe essere il caso che i personaggi che arrivano al livello 4 o 6 o altro vadano in pensione; la campagna si concentra sulle avventure di basso livello. Oppure, potrebbe esserci un limite morbido, con i giocatori incoraggiati a lasciar riposare i loro personaggi migliori mentre continuano a giocare gli altri. Il punto di crisi è evitato perché non si avanza mai abbastanza da sbatterci contro.

\paragraph{Smettere di giocare}: dopo aver giocato per un anno o due, siamo riusciti a portare un po' di personaggi fino ai livelli intermedi. Però è bastata una mossa sbagliata e ora il gruppo è morto. Non è uno spazio naturale per porre fine alla campagna? Non so se sia una vera soluzione, però è una fine.

\paragraph{Inflazione del successo}: rendere il successo più facile muove naturalmente l'intero concetto di mid-tier più in là dal punto di vista numerico; puoi giocare un personaggio di sesto livello e sentire come se fosse di terzo secondo la stima precedente. Se è facile arrivare al sesto livello, allora forse è più facile sopportare l'occasionale perdita di un personaggio di sesto livello ogni tanto. Basta farne un altro! Ovviamente, questa è difficilmente una vera soluzione; stiamo semplicemente riorganizzando le sedie sul ponte di una nave che affonda.

\paragraph{Diventare teatrali}: storicamente, la linea principale di sviluppo di \dnd{} ha risolto il punto di crisi passando da una sfida vera a una sfida finta. Aggiungere abbastanza bilanciamento degli incontri, regole di resurrezione e imbrogli coi dadi porterà indubbiamente ad avere una partita di \dnd{} che sia abbastanza sicura da poterla giocare senza doversi preoccupare troppo di perdere più personaggi di quanti il gruppo non possa permettersi dal punto di vista emotivo.

\paragraph{Nuove modalità di gioco}: il gioco di dominio ai livelli nominati non è mai diventato un qualcosa di completamente stabilito in \dnd{}, ma ha suggerito una soluzione in stile \wargam{e} al punto di crisi. Se il gioco cambia ai livelli alti, cambiando tematiche al punto che la morte del personaggio non è l'esatta definizione di sconfitta, allora possiamo ancora avere un \wargam{e} anche con personaggi troppo importanti per rischiarli regolarmente, giusto?

Parlo del punto di crisi come qualcuno che lo vede da sotto, ma non sa davvero cosa vi si trovi oltre. Forse qualcun altro ha giocato una partita di \dnd{} difficile, legittimamente sfidante fino alle altezze a cui perdere il proprio personaggio diventa inaccettabile e può raccontarci se esiste un simile gioco là fuori. I spero di poter esplorare queste possibilità da me, non appena riusciremo ad avanzare la nostra campagna fino al punto di crisi.